\clearpage
\section{Implementation map, execution order, and venue strategy}

\subsection{Current versus proposed code}

\begin{longtable}{P{0.19\textwidth}P{0.37\textwidth}P{0.36\textwidth}}
\caption{Implementation map. Remaining proposed items must not be described as implemented.}\label{tab:module-map}\\
\toprule
Phase & Current reusable code & Proposed implementation \\
\midrule
\endfirsthead
\toprule
Phase & Current reusable code & Proposed implementation \\
\midrule
\endhead
0 & SAFE schema/dataset validators; cluster runbook; frozen stage runner. & Experiment contract, allocation validator, and cross-manifest identity audit. \\
1 & \path{run_stage_aware_parent_safe.py}; \path{stage_aware_parent_safe.py}; frozen documentation and tests. & Confirmation count/support validator; CI criterion in the frozen gate; deliberate CI-failure test. \\
2 & \path{full_snapshot.py}; \path{counterfactual_branch.py}; Xiaomi policy \code{get\_state/set\_state}; exact replay runner, printer, combined audit, and focused tests. & Conditional task/platform extensions only; the bounded two-task gate has passed. \\
3 & Phase~2 candidate payloads; frozen terminal SAFE scorer; ordered subtask evaluator; \path{analyze_phase3_candidate_opportunity.py} and tests. & Conditional five-task collection plan, support validator, and opportunity-video sampler after the bounded pilot gate. \\
4 & Recovery modes, ordered progress/regression trace, policy-pluggable benchmark. & Operator API, delay-grid collector, common-randomness provenance, and recoverability-surface analysis. \\
5 & Raw action-token features; proprioceptive state history; stage-aware weighting ideas. & Candidate-specific critic dataset/model/runner and causal stage tracker. \\
6 & Deterministic random/lowest/highest candidate selection and mock tests. & Recovery-value scorer interface, sealed-snapshot arm runner, and matched candidate-set audit. \\
7 & Genuine-inference feature records and post-failure recovery hooks. & Pre-action Xiaomi monitor, utility gate/router, intervention budget, and online randomized evaluator. \\
8 & Detached pre-output hidden feature and client-side no-gradient scoring. & New model/server patch for noise-step critic and differentiable flow guidance. \\
9 & RLDX raw feature adapter and separate official recovery path. & Second-policy snapshot/session restoration, frozen shift wrappers, and external physical safety infrastructure. \\
\bottomrule
\end{longtable}

The current Xiaomi ``observation context'' is flattened proprioceptive state history, not a visual embedding. The report and code should call it \emph{state context}. A visual-context ablation requires explicit causal RGB embeddings and their own storage, normalization, and domain-shift evaluation.

\subsection{Execution order}

Phase 1 is complete and failed its frozen confirmation gate. Phase 2 is complete and passed all bounded exact-replay gates. Phase 3 is now active, beginning with a read-only opportunity analysis of the validated two-task branches. No simulator server or new collection is required for this first step. Critic training must not begin until Phase~3 establishes causal candidate/operator headroom; when it does begin, validation and test parents remain sealed.

The following dependencies must remain sequential:
\begin{enumerate}
  \item exact snapshot replay before counterfactual scale;
  \item oracle opportunity before critic training;
  \item sealed offline critic gate before prospective reranking;
  \item prospective reranking before differentiable guidance;
  \item development freeze before every fresh confirmation or final trial;
  \item healthy ordinary-policy positive control before second-policy failure collection.
\end{enumerate}

\subsection{Cluster and artifact discipline}

Source code lives under \path{/gs/fs/tga-shinoda/felid}; environments, checkpoints, datasets, rollouts, videos, logs, and reports live under \path{/gs/bs/tga-shinoda/felid}. Start and verify the policy server before the simulator client. Use \code{python -u} under background execution, unique ports and output roots, disabled W\&B by default, immutable manifests, and persistent logs. Storage quota or inode errors are stop-and-inspect events. Preserve partial evidence and resume or merge rather than deleting failed blocks.

\subsection{Venue completion bars}

\begin{table}[H]
\centering
\caption{Recommended evidence packages. These are competitive research recommendations, not official acceptance rules.}
\label{tab:venues}
\small
\begin{tabularx}{\textwidth}{P{0.13\textwidth}Y P{0.28\textwidth}}
\toprule
Venue & Competitive package & Main risk \\
\midrule
CoRL & Causal benchmark; learned utility gate; held-out tasks/scenes; prospective closed-loop gain; preferably second policy or physical evidence. & Five-task detector or oracle-heavy recovery appears incremental beside CheckVLA and B2FF. \\
RSS & Exceptionally clean exact-state causal design, deadline/latency and safety analysis, strong physical-system evidence, and direct comparison to the closest recovery systems. & Broad benchmark tables without a crisp robotics mechanism. \\
ICRA & Complete integrated system, broad RoboCasa365 simulation, runtime/safety analysis, strong component ablations, and several physical multi-stage tasks where feasible. & Simulator-only one-backbone evaluation may look incomplete. \\
CVPR/ICCV & Recoverability as a visual/action representation contribution, multiple encoders, held-out camera/object/scene shifts, and closed-loop control validating the vision method. & A primarily trigger/control paper is a scope mismatch. \\
WACV & Reproducible causal benchmark/application, public code/data, rigorous negative audit, critic reranking, and meaningful visual/context evaluation. & Weak fit without released artifacts or significant vision/application content. \\
\bottomrule
\end{tabularx}
\end{table}

\subsection{Recommended paper contributions}

\begin{enumerate}
  \item We introduce intervention-specific recoverability curves for long-horizon VLA execution, measured through exact-prefix counterfactual branches over task stages, delays, and recovery operators.
  \item We show that terminal failure likelihood is not equivalent to causal intervention value or remaining recovery time.
  \item We learn a stage-aware utility gate and candidate critic that improve prospective task success while controlling unnecessary and harmful intervention.
\end{enumerate}

Critic-guided flow is a fourth contribution only if Phase 8 passes matched-compute and held-out safety tests. Otherwise, discrete critic reranking is the final executor.

\subsection{Immediate next actions}

\begin{enumerate}
  \item Preserve the failed Phase~1 bundles and the passing Phase~2 root, plan, branch payloads, hashes, and error ledger as immutable evidence.
  \item Run \path{analyze_phase3_candidate_opportunity.py} on the existing 20 snapshots and 80 candidate branches; optionally score the identical sets with the frozen independent SAFE ensemble.
  \item Inspect task/stage support, mixed outcomes, oracle headroom, successful-control harm, numerical ties, and hierarchical intervals without changing any Phase~2 artifact.
  \item If either bounded pilot criterion fails, stop critic work and predeclare one redesign of proposal diversity, timing, horizon, or operator.
  \item If both estimable criteria pass, freeze and collect the five-task Phase~3 opportunity screen needed to evaluate the formal three-of-five-task gate.
  \item Scale recovery surfaces only if the formal oracle-headroom and support gates pass.
  \item Train simple recovery-value probes before any broad neural sweep.
  \item Run the 200-rollout fixed-trigger critic-selection pilot.
  \item Power the primary Phase 7 comparison, then consider flow guidance and external validation.
\end{enumerate}

\cautionnote{First research priority.}{Do not launch another detector sweep, train a candidate critic, or scale to the approximately 1,000-suffix formal screen first. Exact replay has passed. The highest-value next evidence is the inexpensive, frozen two-task oracle-opportunity analysis of branches already collected.}
